\section{第二节
Java/Python基础与应用}\label{ux7b2cux4e8cux8282-javapythonux57faux7840ux4e0eux5e94ux7528}

\subsection{本节内容}\label{ux672cux8282ux5185ux5bb9}

本节将介绍Java和Python两种主流后端开发语言在智慧水利平台中的应用，包括基础语法、核心特性、开发框架以及实际应用案例等内容。

\subsubsection{本节目录}\label{ux672cux8282ux76eeux5f55}

\begin{itemize}
\tightlist
\item
  \href{section05-02/section05-02-01.md}{5.2.1
  Java与Python在水利领域的应用概述}
\item
  \href{section05-02/section05-02-02.md}{5.2.2 Java基础与应用}
\item
  \href{section05-02/section05-02-03.md}{5.2.3 Python基础与应用}
\item
  \href{section05-02/section05-02-04.md}{5.2.4 Java与Python的协同使用}
\item
  \href{section05-02/section05-02-05.md}{5.2.5 智慧水利平台案例分析}
\end{itemize}

\subsection{学习要点}\label{ux5b66ux4e60ux8981ux70b9}

\begin{itemize}
\tightlist
\item
  理解Java和Python在水利信息化中的应用场景和优势
\item
  掌握Java和Python的基础语法和核心特性
\item
  了解主流框架在智慧水利平台开发中的应用
\item
  熟悉两种语言的协同使用策略和方法
\item
  通过案例分析理解语言选择和技术架构设计
\end{itemize}

\subsection{思考题}\label{ux601dux8003ux9898}

\begin{enumerate}
\def\labelenumi{\arabic{enumi}.}
\tightlist
\item
  比较Java和Python在智慧水利平台开发中的优缺点，并分析它们各自适合的应用场景。
\item
  在一个大型水利工程监测系统中，如何合理规划Java和Python的应用边界？
\item
  设计一个基于Java的水文数据管理API，要求支持数据的增删改查及简单统计。
\item
  使用Python开发一个简单的洪水预警模型，考虑输入的水文数据和输出的预警结果。
\item
  探讨在微服务架构下，Java服务和Python服务如何有效协作。
\end{enumerate}

\subsection{参考文献}\label{ux53c2ux8003ux6587ux732e}

\begin{enumerate}
\def\labelenumi{\arabic{enumi}.}
\tightlist
\item
  Joshua Bloch. ``Effective Java, 3rd Edition''. Addison-Wesley, 2018.
\item
  Mark Lutz. ``Learning Python, 5th Edition''. O'Reilly Media, 2013.
\item
  Craig Walls. ``Spring Boot in Action''. Manning Publications, 2016.
\item
  Jake VanderPlas. ``Python Data Science Handbook''. O'Reilly Media,
  2016.
\item
  Sam Newman. ``Building Microservices''. O'Reilly Media, 2015.
\end{enumerate}
